\section{Discussion and Limitations}
\label{sec:disc}

C1 (Proposition~\ref{prop:retain}) is conditional true-fault retention under
hypotheses R1--R5, not a competitor to full-history timed
diagnosis~\cite{koehl2023}. Monotonicity (Proposition~\ref{prop:mono}) is a
separate statement. C2 (Proposition~\ref{prop:c2}) is a one-step operator on
\((s(t),\mathrm{cost}(t),\mathrm{id})\), not the closed-loop supervisors or
conditional plans already in the
literature~\cite{sampath1998,chanthery2009,hu2020,miao2025}. Constraint-solving
test generation remains a neighbouring generator, not this paper's selection
rule~\cite{petrenko2016}. Aerospace procedure work remains the engineering
setting and a localization-risk warning~\cite{chanthery2019}.

\phbox{PH-RES-INTERPRET}{EXPERIMENT-RESULT}{Discussion of measured gains or
losses waits for the DETECT, LOCATE and ABLATION cells.}

\paragraph{Limitations.}
The observation model is conservative: missing, reordered or time-skewed
events are not reconstructed, so C1 does not extend to a general loss, drift
or reordering theorem. R1--R5, exclusive ownership and the exclusive resource
mode are hypotheses of the specification, not measured platform properties.
Runtime ARINC~645 capability is not established. Equivalence among mutants and
among protocol traces can leave distinguishable faults unresolved. Clock and
error budgets remain publication decisions~\cite{jcgm100,jcgm106}.
Confirmatory measurements have not yet been obtained; this manuscript does
not claim measured improvement, implementation completeness, or certification.
